\documentclass[11pt,a4paper,openany]{book}
\input{userguide-preamble}
\hypersetup{pdftitle={HugMesh Games --- Godot developer walkthrough}}
\begin{document}

\begin{titlepage}\thispagestyle{empty}
\noindent\begin{tikzpicture}[remember picture,overlay]
  \fill[snapBg] (current page.south west) rectangle (current page.north east);
  \fill[snapAccent,opacity=0.18]    ([yshift=-0.05\paperheight]current page.north west) circle (8cm);
  \fill[snapAccentTwo,opacity=0.10] ([xshift=4cm,yshift=-2cm]current page.north west) circle (5cm);
\end{tikzpicture}
\vspace*{4cm}\noindent
{\sffamily\bfseries
  {\color{snapAccentTwo}\fontsize{18}{22}\selectfont HUGMESH GAMES}\\[3mm]
  {\color{white}\fontsize{42}{50}\selectfont Godot\\developer guide}\par}
\vspace{1.5cm}\noindent\color{white!75!snapAccentTwo}\rule{0.4\linewidth}{1.2pt}
\vspace{0.8cm}\noindent
{\sffamily\color{white!85}\large
From \texttt{mkdir games/my-game/} to a row on the home page: a
step-by-step Godot 4 walkthrough.}
\vfill\noindent
{\sffamily\color{white!70}\small
Companion to \texttt{godot-developer.pptx}. Generated: \snapbuilddate{} \snapbuildtime}
\end{titlepage}

\frontmatter\pagestyle{empty}\tableofcontents\clearpage
\mainmatter\pagestyle{fancy}

\chapter{What you'll build}
A new Godot 4.6 project under \snapfile{games/<game-slug>/} that opens
standalone in the editor. It uses the reusable Rust GDExtension at
\snapfile{extensions/hugmesh-bridge/} for HostBridge signals. The
HugMesh Export plugin produces a deploy-ready zip with
\texttt{manifest.json} (schema v2) + \texttt{register.sql} +
\texttt{bundle/}. You ship it either through \texttt{make build-godot}
(committed in-repo) or through \texttt{/admin/games} (hot upload).

\chapter{Prerequisites}
\begin{itemize}
  \item Godot 4.6 in \texttt{PATH} as \texttt{godot}.
  \item Rust stable (\texttt{cargo --version}).
  \item Optional: Emscripten + rustup \texttt{wasm32-unknown-emscripten}
        for the Rust GDExtension WASM target.
\end{itemize}
Run once: \texttt{make build-rust} so every \snapfile{games/<dir>/bin/}
gets a fresh \texttt{libhugmesh\_bridge.so}.

\chapter{Step 1 --- Create the project}
\begin{enumerate}
  \item \texttt{mkdir games/my-game-slug/}
  \item Open Godot $\rightarrow$ New Project $\rightarrow$ set path
        $\rightarrow$ Create.
  \item Copy \snapfile{hugmesh\_bridge.gdextension} and
        \snapfile{addons/hugmesh\_export/} from
        \snapfile{games/three-wizards-in-a-trenchcoat/}.
  \item Drop a starter scene at \snapfile{scenes/main.tscn}.
\end{enumerate}

\chapter{Step 2 --- Pick your slugs}
\textbf{GAME slug} appears in \texttt{games.slug},
\texttt{federation\_modules.slug}, and (sometimes) the directory name
under \snapfile{games/}. \textbf{LEVEL slug(s)} --- one per export
preset; appear in \texttt{levels.slug}, the export path, and the
Makefile.

Lowercase, hyphen-free, kebab-case (\texttt{forest-glade}, not
\texttt{ForestGlade}). Once committed, slugs do not change without a
coordinated DB migration. Pick well.

\chapter{Step 3 --- Wire HostBridge}
The reusable GDExtension exposes a \texttt{HostBridge} node with typed
signals: \texttt{movement(x,y,z,action,pressed)},
\texttt{world\_state(state)}, \texttt{character\_state(state)},
\texttt{scene\_state(state)}, \texttt{time(state)}.

Add a HostBridge node to the scene tree, then connect signals from your
gameplay scripts. See
\snapfile{games/three-wizards-in-a-trenchcoat/scripts/scene\_loader.gd}
for a reference subscriber.

\chapter{Step 4 --- Author export presets}
Project $\rightarrow$ Export $\rightarrow$ Add new preset $\rightarrow$
Web $\rightarrow$ name it after the level slug (\texttt{ForestGlade}).

\begin{lstlisting}[language=bash]
export_path  = ../../backend/src/main/resources/assets/levels/forest-glade/index.html
custom_features = level_forest-glade
exclude_filter  = *.gdextension
\end{lstlisting}

Save the \texttt{.cfg}, do not export from the editor. The Makefile
drives headless export.

\chapter{Step 5a --- Commit and ship}
\begin{enumerate}
  \item Add \snapfile{games/<dir>/} to \texttt{GAME\_PROJECTS} in the
        root \texttt{Makefile}.
  \item Define \texttt{PRESETS\_<dir>} with the preset names.
  \item Write a Liquibase changeset for \texttt{games},
        \texttt{levels}, \texttt{federation\_modules},
        \texttt{game\_access} (use
        \snapfile{017-snapdragon-catalog.xml} as the template).
  \item Commit. CI runs \texttt{make build} and your game lands.
\end{enumerate}

\chapter{Step 5b --- Upload through admin}
\snapdiagram{game-upload-controller.tex}

In Godot: \emph{Project $\rightarrow$ Tools $\rightarrow$ HugMesh
$\rightarrow$ Export\ldots}. Fill in slug, display name, description,
scene type. The plugin writes a deploy-ready zip with the v2
\texttt{manifest.json} + \texttt{register.sql} + \texttt{bundle/}.

Visit \texttt{/admin/games} on the live stack and drag the zip into the
upload card. The BFF validates manifest schema v2, hashes the bundle,
atomically moves to \texttt{<gamesDir>/<slug>/<sha-prefix>}, and runs
\texttt{register.sql} in one R2DBC transaction.

\begin{figure}[!h]\centering
\includegraphics[width=0.92\linewidth]{admin-uploads.png}
\caption{The admin upload card with rollback history below.}
\end{figure}

\chapter{Local test loop}
\begin{enumerate}
  \item \texttt{docker compose up -d postgres cassandra casdoor api wasm-fragment}.
  \item In Godot, press Play. \snapfile{scripts/api\_client.gd} auto-targets
        \texttt{http://localhost:8080}.
  \item \snapfile{scripts/scene\_loader.gd} polls
        \texttt{/api/levels/\{id\}/scene} and spawns elements.
  \item Edit a scene; press Play again. No build step required for
        local iteration.
  \item Ready to ship: \texttt{make build-godot \&\& docker compose
        restart api wasm-fragment}.
\end{enumerate}

\chapter{Multiplayer-aware bridge}
The Rust HostBridge gained two methods the single-player demo didn't
need:

\begin{itemize}
  \item \texttt{register\_player\_slot(slot, character\_id)} ---
        called from \snapfile{games/<slug>/scripts/multiplayer\_lobby.gd}
        once per player as they join. The GDScript reference is
        \snapfile{multiplayer\_lobby.gd}; it spawns one wizard
        \texttt{CharacterBody3D} per registered slot.
  \item \texttt{view\_egress(camera\_pos, look\_dir)} --- $\sim$10\,Hz
        posts to \texttt{/api/view}. The BFF fans these out on the
        per-session \emph{view} sink that the admin console at
        \texttt{/admin/sessions/\{id\}/stream} subscribes to. Don't
        rate-limit yourself; the nginx \texttt{sse} zone caps at
        30 r/s per IP anyway.
\end{itemize}

Two contracts your scene must respect:

\begin{enumerate}
  \item Every \texttt{POST /api/controls/\{character\_id\}} must
        carry an \texttt{X-Session-Id} header. Without it, two players
        sharing the same character record across sessions race
        Postgres' first-row pick at
        \texttt{findActiveOwnerOfCharacter(character\_id)} and one
        side gets 403. The \snapfile{game\_bridge.gd} api client
        already adds the header --- if you fork it, keep that line.
  \item On a 403 from \texttt{/api/controls}, post
        \texttt{SD\_SESSION\_EXPIRED} to the parent window. The
        Angular wrapper's \snapfile{ApiBridgeService.handleSessionExpired()}
        will re-run the session.connect $\rightarrow$ claimControl
        chain automatically.
\end{enumerate}

The HostBridge build is the same one used by the RPG + VN makers ---
those bundles already ship \snapfile{hugmesh\_bridge.gdextension}
in the produced source tree, so multiplayer-awareness comes for free
when a publisher uses a maker rather than authoring directly in Godot.

\chapter{Two-player local test}
The browser-side multiplayer contract is pinned by
\reposrc{e2e/features/multiplayer/rpg-two-players.feature}: two real
Casdoor users open the lobby, join the same session, launch the same
chapter; each one's control event must reach the other's iframe via
the per-session SSE sink.

To exercise it locally without Playwright:

\begin{enumerate}
  \item Open two browser profiles (e.g. one normal + one incognito).
        \emph{Cookies must be isolated} --- shared cookies race the
        session tracker.
  \item In each, log in as a different player (the
        \texttt{player1} and \texttt{player2} seed users share one
        password on the integration-test profile — a fresh random
        value \texttt{make e2e} generates and writes to
        \texttt{.e2e-credentials}; export \texttt{E2E\_PASSWORD=...}
        yourself instead if you want to pin a known value).
  \item Profile A: open the game, \emph{Create lobby}, public
        visibility. Profile B: open the same game, \emph{Join} on the
        new row.
  \item Both click \emph{Launch}. The WASM compile is 15--90\,s on a
        cold image --- gate input on the \texttt{?sessionId=\ldots}
        EventSource open before pressing keys, or the spec races
        the subscription.
  \item Move profile A; profile B's iframe should see the movement
        within a few hundred ms.
\end{enumerate}

The four canvas-window controls (Fullscreen, Open-in-new-tab,
Open-in-new-window with pause, Resume) are covered by
\reposrc{e2e/features/multiplayer/window-controls.feature}. The
``open in new window'' button posts \texttt{SD\_GAME\_PAUSE} to the
parent iframe so only one window plays at a time --- if you wire a
custom toolbar, mirror that behaviour.

\chapter{Common mistakes}
\begin{itemize}
  \item Slugs differ across the four contract points $\rightarrow$
        iframe 404s or wrong scene loads.
  \item \texttt{exclude\_filter} misses \texttt{*.gdextension}
        $\rightarrow$ bundle balloons by 8 MB.
  \item Calling \texttt{add\_child} from \texttt{\_ready} on a node
        already in the tree $\rightarrow$ ``parent is busy''. Use
        \texttt{call\_deferred}.
  \item Forgetting to run \texttt{make build-rust} after pulling
        $\rightarrow$ stale \snapfile{games/*/bin/.so}.
\end{itemize}

\end{document}
